\documentclass[12pt, a4paper]{article}
\usepackage[UTF8]{ctex}
\usepackage{amsmath, amssymb, amsthm}
\usepackage{geometry}
\usepackage{bm}

\geometry{left=2.5cm, right=2.5cm, top=2.5cm, bottom=2.5cm}

\newcommand{\RR}{\mathbb{R}}
\newcommand{\e}{\mathrm{e}}
\newcommand{\dif}{\mathrm{d}}

\begin{document}
\section*{14.1}
\textbf{1.}求下列第一类曲线积分：

(2)\(\int\limits_L |y| \dif s\)，其中\(L\)为单位圆周\(x^2 + y^2 = 1\).

\[\int\limits_L |y| \dif s = \int_{0}^{2\pi} |\sin t| \dif t = 4.\]

(4)\(\int\limits_L |x| \dif s\)，其中\(L\)为双纽线\((x^2 + y^2)^2 = x^2 - y^2\).

\[\int\limits_L |x| \dif s = \int_{-\frac{\pi}{4}}^{\frac{\pi}{4}} |\cos\theta| \dif \theta + \int_{\frac{3\pi}{4}}^{\frac{5\pi}{4}} |\cos\theta| \dif \theta = 2\sqrt{2}.\]

(5)\(\int\limits_L (x^2 + y^2 + z^2) \dif s\)，\(L\)为一段螺旋线\(x = a\cos t, y = a\sin t, z = bt, 0 \leq t \leq 2\pi\).

\[\int\limits_L (x^2 + y^2 + z^2) \dif s = \sqrt{a^2 + b^2}\left(2\pi a^2 + \frac{8\pi^3}{3}b^2\right).\]

(7)\(\int\limits_L (xy + yz + zx) \dif s\)，其中\(L\)为球面\(x^2 + y^2 + z^2 = a^2\)和平面\(x + y + z = 0\)的交线.

\[\int\limits_L (xy + yz + zx) \dif s = \int\limits_L \frac{1}{2}\left[(x + y + z)^2 - (x^2 + y^2 + z^2)\right] \int\limits_L \frac{-a^2}{2} \dif s\dif s = -\pi a^3.\]

\textbf{3.}求下列曲面的面积：

(3)球面\(x^2 + y^2 + z^2 = a^2\)包含在锥面\(z = \sqrt{x^2 + y^2}\)内的部分.

\[x = a\sin\varphi\cos\theta, y = a\sin\varphi\sin\theta, z = a\cos\varphi, \quad \varphi \in [0, \pi], \theta \in [0. 2\pi].\]
\[z \geq \sqrt{x^2 + y^2} \Rightarrow \varphi \in \left[0, \frac{\pi}{4}\right].\]
\[S = \int_{0}^{2\pi} \dif \theta \int_{0}^{\frac{\pi}{4} a^2\sin\varphi \dif \varphi = \pi a^2(2 - \sqrt{2})}.\]

(5)抛物面\(x^2 + y^2 = 2az\)包含在柱面\((x^2 + y^2)^2 = 2a^2xy\)内的部分.

\[x = r\cos\theta, y = r\sin\theta, r^2 = a^2\sin 2\theta > 0, \theta \in \left[0, \frac{\pi}{2}\right] \cup \left[\pi, \frac{3\pi}{2}\right].\]
\[\dif S = \sqrt{1 + z_x^2 + z_y^2} \dif x \dif y = \frac{1}{a}\sqrt{a^2 + r^2}r \dif r \dif \theta.\]
\[S = 2\int_{0}^{\frac{\pi}{2}} \dif \theta \int_{0}^{a\sqrt{\sin 2\theta}} \frac{1}{a}\sqrt{a^2 + r^2} \dif r = \frac{a^2}{9}(20 - 3\pi)\]

(6)环面\(\begin{cases}
x = (b + a\cos\theta)\cos\varphi, \\
y = (b + a\cos\theta)\sin\varphi, \\
z = a\sin\theta,
\end{cases} 0 \leq \theta \leq 2\pi, 0 \leq \varphi \leq 2\pi\)，其中\(0 < a < b\).

\[\vec{r}_\theta = (-a\sin\theta\cos\varphi, -a\sin\theta\sin\varphi, a\cos\theta),\]
\[\vec{r}_\varphi = (-(b + a\cos\theta)\sin\varphi, (b + a\cos\theta)\cos\varphi, 0).\]
\[\vec{r}_\theta \times \vec{r}_\varphi = -a(b + a\cos\theta)(\cos\theta\cos\varphi, \cos\theta\sin\varphi, \sin\theta),\]
\[\|\vec{r}_\theta \times \vec{r}_\varphi\| = a(b + a\cos\theta).\]
\[S = \int_{0}^{2\pi} \dif \varphi \int_{0}^{2\pi} a(b + a\cos\theta) \dif \theta = 4\pi^2ab.\]

\textbf{4.}求下列第一类曲面积分：

(2)\(\iint\limits_\Sigma (x^2 + y^2) \dif S\)，其中\(\Sigma\)是区域\(\{(x, y, z)|\sqrt{x^2 + y^2} \leq z \leq 1\}\)的边界.

\begin{align*}
\iint\limits_\Sigma (x^2 + y^2) \dif S &= \iint\limits_{\substack{z = 1 \\ x^2 + y^2 \leq 1}} (x^2 + y^2) \dif S + \iint\limits_{\substack{z = \sqrt{x^2 + y^2} \\ z \leq 1}} (x^2 + y^2) \dif S \\
&= \int_{0}^{2\pi} \dif \theta \int_0^1 r^3 \dif r + \sqrt{2}\iint\limits_D (x^2 + y^2) \dif x \dif y \\
&= \frac{\pi}{2}(1 + \sqrt{2}).
\end{align*}

(5)\(\iint\limits_\Sigma \left(\dfrac{x^2}{2} + \dfrac{y^2}{3} + \dfrac{z^2}{4} \dif S\right)\)，其中\(\Sigma\)是球面\(x^2 + y^2 + z^2 = a^2\).

\[\iint\limits_\Sigma x^2 \dif S = \iint\limits_\Sigma y^2 \dif S = \iint\limits_\Sigma z^2 \dif S = \frac{1}{3} \iint\limits_\Sigma (x^2 + y^2 + z^2) \dif S = \frac{4\pi a^4}{3}.\]
\[\iint\limits_\Sigma \left(\dfrac{x^2}{2} + \dfrac{y^2}{3} + \dfrac{z^2}{4} \dif S\right) = \frac{13\pi a^4}{9}.\]

(7)\(\iint\limits_\Sigma z \dif S\)，其中\(\Sigma\)是螺旋面\(x = u\cos v, y = u\sin v, z = v, 0 \leq u \leq a, 0 \leq v \leq 2\pi\)的一部分.

\[\vec{r}_u = (\cos v, \sin v, 0), \vec{r}_v = (- u\sin v, u\cos v, 1).\]
\[\vec{u} \times \vec{r}_v = (\sin v, -\cos v, u), \left\|\vec{u} \times \vec{r}_v\right\| = \sqrt{1 + u^2}.\]
\[\iint\limits_\Sigma z \dif S = \int_{0}^{2\pi} \dif v \int_0^a v\sqrt{1 + u^2} \dif u = \pi^2\left[a\sqrt{1 + a^2} + \ln\left(a + \sqrt{1 + a^2}\right)\right].\]

\textbf{8.}设\(u(x, y ,z)\)为连续函数，它在\(M(x_0, y_0, z_0)\)处有连续的二阶偏导数.记\(\Sigma\)为以\(M\)点为中心，半径为\(R\)的球面，以及
\[T(R) = \frac{1}{4\pi R^2}\iint\limits_\Sigma u(x, y, z) \dif S.\]

(1)证明：\(\lim\limits_{R \to 0} T(R) = u(x_0, y_0, z_0)\).

(2)若\(\left.\left(\dfrac{\partial^2 u}{\partial x^2} + \dfrac{\partial^2 u}{\partial y^2} + \dfrac{\partial^2 u}{\partial z^2}\right)\right|_{(x_0, y_0, z_0)} \neq 0\)，求当\(R \to 0\)时无穷小量\(T(R) - u(x_0, y_0, z_0)\)的主要部分.

解：(1)\(\forall \varepsilon > 0, \exists \delta > 0\)，当\(\sqrt{(x - x_0)^2 + (y - y_0)^2 + (z - z_0)^2} < \delta\)时，
\[|u(x, y, z) - u(x_0, y_0, z_0)| < \varepsilon.\]
当\(R < \delta\)时，
\begin{align*}
|T(R) - u(x_0, y_0, z_0)| &= \left|\frac{1}{4\pi R^2}\iint\limits_\Sigma u \dif S - u(x_0, y_0, z_0)\right| \\
&= \left|\frac{1}{4\pi R^2}\iint\limits_\Sigma [u - u(x_0, y_0, z_0)] \dif S\right| \\
&\leq \frac{1}{4\pi R^2}\iint\limits_\Sigma |u - u(x_0, y_0, z_0)| \dif S \\
&< \varepsilon.
\end{align*}

(2)令
\[\begin{cases}
x = x_0 + R\xi, \\
y = y_0 + R\eta, \\
z = z_0 + R\zeta,
\end{cases}\]
\begin{align*}
u(x, y, z) = &u(x_0, y_0, z_0) + R\left(\xi\frac{\partial}{\partial x} + \eta\frac{\partial}{\partial y} + \zeta\frac{\partial}{\partial z}\right)u(x_0, y_0, z_0) \\
&+ \frac{R^2}{2}\left(\xi\frac{\partial}{\partial x} + \eta\frac{\partial}{\partial y} + \zeta\frac{\partial}{\partial z}\right)^2u(x_0, y_0, z_0) + o(R^2).
\end{align*}
\begin{align*}
T(R) - u(x_0, y_0, z_0) &= \frac{1}{4\pi R^2}\iint\limits_\Sigma [u - u(x_0, y_0, z_0)] \dif S \\
&= \frac{R}{4\pi} \iint\limits_\sigma \left(\xi\frac{\partial}{\partial x} + \eta\frac{\partial}{\partial y} + \zeta\frac{\partial}{\partial z}\right)u(x_0, y_0, z_0) \dif \sigma \\
&+ \frac{R^2}{8\pi} \iint\limits_\sigma \left(\xi\frac{\partial}{\partial x} + \eta\frac{\partial}{\partial y} + \zeta\frac{\partial}{\partial z}\right)^2u(x_0, y_0, z_0) \dif \sigma \\
&= \frac{R^2}{6}\left.\left(\dfrac{\partial^2 u}{\partial x^2} + \dfrac{\partial^2 u}{\partial y^2} + \dfrac{\partial^2 u}{\partial z^2}\right)\right|_{(x_0, y_0, z_0)}
\end{align*}

\textbf{9.}设\(\Sigma\)为上半椭球面\(\dfrac{x^2}{2} + \dfrac{y^2}{2} + z^2 = 1(z \geq 0)\)，\(\pi\)为\(\Sigma\)在点\(P(x, y, z)\)处的切平面，\(\rho(x, y, z)\)为原点\(O(0, 0, 0)\)到平面\(\pi\)的距离，求\(\iint\limits_\Sigma \dfrac{z}{\rho(x, y, z)} \dif S\).

解：
\[\pi:xX + yY + 2zZ = 2.\]
\[\rho = \frac{2}{\sqrt{x^2 + y^2 + 4z^2}}.\]
\begin{align*}
\iint\limits_\Sigma \dfrac{z}{\rho(x, y, z)} \dif S &= \iint\limits_\Sigma \dfrac{z\sqrt{x^2 + y^2 + 4z^2}}{2} \dif S \\
&= \iint\limits_\Sigma \dfrac{z\sqrt{x^2 + y^2 + 4z^2}}{2} \frac{\sqrt{x^2 + y^2 + 4z^2}}{2z} \dif x \dif y \\
&= \frac{1}{4}\iint\limits_D (x^2 + y^2 + 4z^2) \dif x \dif y \\
&= \frac{3\pi}{2}.
\end{align*}

\textbf{10.}设\(\Sigma\)是单位球面\(x^2 + y^2 + z^2 = 1\).证明
\[\iint\limits_\Sigma f(ax + by + cz) \dif S = 2\pi \int_{-1}^{1} f(u\sqrt{a^2 + b^2 + c^2}) \dif u,\]
其中\(a, b, c\)为不全为零的常数，\(f(u)\)是\(|u| \leq \sqrt{a^2 + b^2 + c^2}\)上的一元连续函数.

解：令\(\vec{n} = (a, b, c), ax + by + cz = \vec{n} \cdot (x, y, z) = \left\|\vec{n}\right\|\cos\theta\)，其中\(\theta\)是\(\vec{n}\)与\((x, y, z)\)的夹角.
\begin{align*}
\iint\limits_\Sigma f(ax + by + cz) \dif S &= \int_{0}^{2\pi} \dif \varphi \int_0^\pi f(\left\|\vec{n}\right\|\cos\theta)\sin\theta \dif \theta \\
&= 2\pi \int_{-1}^{1} f(u\sqrt{a^2 + b^2 + c^2}) \dif u \quad (u = \cos\theta).
\end{align*}

\section*{14.2}
\textbf{1.}求下列第二类曲线积分：

(2)\(\int\limits_L (x^2 - 2xy) \dif x + (y^2 - 2xy) \dif y\)，其中\(L\)是抛物线的一段：\(y = x^2, -1 \leq x \leq 1\)，方向由\((-1, 1)\)到\((1, 1)\).

\[\int\limits_L (x^2 - 2xy) \dif x + (y^2 - 2xy) \dif y = \int_{-1}^{1} (x^2 - 2x^3 - 4x^4 + 2x^5) \dif x = -\frac{14}{15}.\]

(3)\(\int\limits_L \dfrac{(x + y) \dif x - (x - y) \dif y}{x^2 + y^2}\)，其中\(L\)是圆周\(x^2 + y^2 = a^2\)，方向为逆时针方向.

\[\int\limits_L \dfrac{(x + y) \dif x - (x - y) \dif y}{x^2 + y^2} = \int_{0}^{2\pi} \frac{(a\cos\theta + a\sin\theta)(-a\sin\theta) - (a\cos\theta - a\sin\theta)(a\cos\theta)}{a^2} \dif \theta = -2\pi.\]

(5)\(\int\limits_L x \dif x + y \dif y + (x + y - 1) \dif z\)，其中\(L\)是从点\((1, 1, 1)\)到点\((2, 3, 4)\)的直线段.

\[\begin{cases}
x = 1 + t, \\
y = 1 + 2t, \\
z = 1 + 3t.
\end{cases}\]
\[\int\limits_L x \dif x + y \dif y + (x + y - 1) \dif z = \int_0^1 (6 + 14t) \dif t = 13.\]

(7)\(\int\limits_L (y - z) \dif x + (z - x) \dif y + (x - y) \dif z\)，\(L\)为圆周\(\begin{cases}
x^2 + y^2 + z^2 = 1, \\
y = x\tan\alpha \left(0 < \alpha < \dfrac{\pi}{2}\right),
\end{cases}\)若从\(x\)轴的正向看去，这个圆周的方向为逆时针方向.

\[\int\limits_L (y - z) \dif x + (z - x) \dif y + (x - y) \dif z = \iint\limits_S (-2, -2, -2) \cdot (\sin\alpha, -\cos\alpha, 0) \dif S = 2\pi(\cos\alpha - \sin\alpha).\]

\textbf{3.}方向依纵轴的负方向，且大小等于作用点的横坐标的平方的力构成一个力场.求一个质量为\(m\)的质点沿抛物线\(y^2 = 1 - x\)从点\((1, 0)\)移到\((0, 1)\)时，场力所做的功.

\[\vec{F}(x, y) = (0, -x^2).\]
\[W - \int\limits_L F_x \dif x + F_y \dif y = -m\int_0^1 (1 - t^2)^2 \dif t = -\frac{8m}{15}.\]

\end{document}